\documentclass{article}
\usepackage[
	margin=2cm
]{geometry}
\usepackage{lyluatex}

\title{Wycliffe Choral Society - Warm-Ups}
\author{Joshua Crunden}
\date{}

\begin{document}
\maketitle
\section{Apple and Pineapple}
\lilypondfile[staffsize=17]{apple-and-pineapple}
\begin{enumerate}
	\item Choir to pat quaver rhythm on chest, alternating right and left, emphasising strong beats.
	\item Stop choir in different parts of the bar, then restart from where stopped.
	\item Swap patting for spoken words "apple and pineapple".
	\item Stop and restart, as before.
	\item Introduce parts starting with the basses - unaccompanied if possible.
	\item Stop and restart, as before.
\end{enumerate}

\section{Fun with scales}
\begin{enumerate}
	\item Scale in unison to a vowel slowly (about 70 bpm).
	\item Scale in two groups, second group starting two notes behind (parallel 3rds).
	\item Four groups, two notes apart (7th chords).
	\item Three groups, three notes apart (parallel fourths with occasional tritone).
	\item Four groups, one note apart (parallel seconds).
	\item Eight groups, one note apart (all the notes!).
\end{enumerate}

\section{Three vs Four vs Five}
\lilypondfile[staffsize=17]{3-4-5}
\begin{enumerate}
	\item Words only, stomp on downbeat (3 and 4).
	\item Divide choir in two (3s and 4s together).
	\item Learn melodies (3 and 4).
	\item Divide choir in two, with melody (3 and 4).
	\item Words only for 5.
	\item Learn melody (5).
	\item Divide choir in two with melody (4 and 5).
	\item Add percussion (click in the number and clap on the word "beats").
	\item All in unison, switch between 3, 4 and 5, with melody and stomp, click and clap.
	\item Split choir in three (3, 4, 5). Just words and percussion.
	\item As above, but with melodies as well.
\end{enumerate}

\section{Sing the Scale, Clap the Number}
The singers sing an ascending major scale in any key, singing the corresponding numbers for the degrees of the scale. Show numbers of fingers on right hand to indicate how far up the scale the singers should sing before going back to the tonic (1). On the left hand, indicate which degree of the scale the singers should also clap on.



\end{document}
